% Generated by scripts/export_feature_tables.py -- do not edit by hand.
\begin{description}
  \item[\texttt{profile\_skill\_count}] Count of canonical skills in the candidate profile query. A richer profile narrows BM25 retrieval and gives the overlap features more to work with; a one-skill profile is close to unrankable on content alone.

  \item[\texttt{profile\_seniority\_level}] Ordinal seniority of the originating posting, used as the candidate's own seniority. Combined with the candidate job's seniority it drives the seniority-match feature; kept standalone too because a Staff-level profile and an Intern-level profile should not be treated as interchangeable even before any matching is computed.

  \item[\texttt{profile\_taxonomy\_coverage}] Number of distinct taxonomy categories (languages, data, cloud, ...) represented in the profile. Separates a deep single-category profile from a broad multi-category one, which changes which job families are plausible matches regardless of raw skill count.

  \item[\texttt{candidate\_job\_popularity\_rank}] Normalized dense rank of how common a candidate job's taxonomy category is among training candidates. Backs the popularity baseline directly and gives the re-ranker a content-agnostic prior for job families that are simply posted more often.

  \item[\texttt{candidate\_job\_trailing\_postings}] Count of same-category postings in the 60 days before the cut-off, among jobs known at training time. Lifetime category popularity is dominated by whatever was common years ago; a trailing window tracks hiring demand as of the query time, which is what a system serving requests at the cut-off actually needs.

  \item[\texttt{bm25\_score}] The stage-1 retrieval score itself, carried into stage 2 as a feature. This is the central mechanic of the two-stage design: the re-ranker is not blind to how strongly lexical retrieval already believed in this candidate, it learns how much to trust that score relative to the structured overlap features.

  \item[\texttt{skill\_overlap\_raw}] Count of profile skills that appear verbatim in the candidate's raw (un-normalized) skill mentions. Deliberately naive: it under-counts true overlap whenever a skill is spelled differently ("js" vs "javascript"), which is what makes the gap against skill\_overlap\_jaccard a measurable demonstration of why taxonomy normalization matters.

  \item[\texttt{skill\_overlap\_jaccard}] Jaccard overlap between the profile's canonical skills and the candidate's canonical skills. The taxonomy-normalized counterpart to skill\_overlap\_raw, and the single strongest engineered signal: it is the one feature unavailable to lexical-only retrieval, since BM25 rewards surface-form overlap rather than canonical-skill overlap.

  \item[\texttt{seniority\_match}] One minus the normalized ordinal distance between profile and candidate seniority. A Senior profile against a Staff posting should rank above the same profile against an Intern posting even when the skill overlap is identical.

  \item[\texttt{location\_match}] Binary indicator that the profile's origin location matches the candidate's location, or either side is Remote. Location is a hard constraint for a large share of real candidates and costs nothing to compute, so it is included even though it is not a taxonomy feature.

  \item[\texttt{query\_candidate\_recency\_gap\_days}] Days between the cut-off and the candidate's posting date. A posting that has been open for months is more likely stale or already filled than one posted last week; this is also the feature that most directly exercises the point-in-time contract, since a leaking implementation would let a future candidate posting produce a negative gap.

\end{description}
